% Standalone appendix document - compiles on its own:
%   latexmk -pdf -outdir=build 02-software-architecture.tex
\documentclass[a4paper,11pt]{article}
\usepackage[english]{babel}
\usepackage{../../appendix-style}
\addbibresource{../../../assets/references.bib}

\title{3.2 Software Architecture}
\author{SW4PRJ4 - Group 3}
\date{\today}

\begin{document}
\maketitle

\begin{versionhistory}
  1.0 & 04-11-2026 & Group 3 & Generation of the TypeScript API client \\
\end{versionhistory}

\section{Overview}
The customer app and the admin panel have no hand-written code for calling the API. They use the API client package instead, and that package is generated from the API's own description of itself. This appendix describes how that works. It covers the chain from the C\# code to the two apps, when we generate the client again, and what happens when the API changes in a way that breaks the apps. Why we chose a generated client is described in appendix 2.1.

\section{From C\# to the Apps}
Figure~\ref{fig:api-client-generation} shows the chain from top to bottom. The top zone is C\# in the Web API, the bottom zone is TypeScript in the API client and the two apps, and NSwag is the generation step in between.

\drawiofig[\linewidth]{api-client-generation}{How the TypeScript API client is generated from the C\# API.}{fig:api-client-generation}

\paragraph{1. Controllers and contracts.} We write the contract in C\# in the Api project. The controllers define the routes, and the contracts define the data in requests and responses. Each action also has a name and states which status codes it can return and what type comes back. All of this ends up in the client.

\paragraph{2. The OpenAPI document.} The controllers do not produce the document themselves, and the document is not a file in the repository. ASP.NET Core's OpenAPI package builds it while the API runs. When someone asks for the document, the package looks at the controllers, their attributes and the contract types, and writes an OpenAPI document that describes every endpoint \cite{microsoft-aspnetcore-openapi}\cite[sec.~3.2.1]{desanctis-web-apis-aspnet-core}. OpenAPI is a standard format for describing an HTTP API, and it does not depend on a programming language \cite{openapi-specification}. The API only serves the document when it runs in Development, so it is not public in production.

\paragraph{3. NSwag.} NSwag reads the document and writes a TypeScript client \cite{nswag-typescript-client-generator}. We run it as a local .NET tool with a fixed version, 14.7.1, so everybody in the group generates exactly the same code. Its settings are in a file in the package. The most important ones are:
\begin{itemize}
  \item The client uses the standard fetch function, which works both in the browser and in React Native.
  \item There is one client class for each controller. The name of an action decides which class it goes in and what the method is called, so the action that lists the items becomes a method on the client for items.
  \item Each client class also gets an interface, so a test can replace it with a fake.
\end{itemize}

\paragraph{4. The generated client.} NSwag writes all the clients and types to one file in the package. Nobody edits that file by hand. If something in the client looks wrong, we fix the C\# code and generate the client again.

\paragraph{5. The client configuration.} One small file in the package is written by hand. It gives the generated clients the fetch function they use. That function adds the user's access token to every call, and if the API answers 401, it tries to renew the token once. Each app tells the package where its token is kept, so the generated code does not need to know anything about Keycloak or about how the app stores the token.

\paragraph{6. The apps.} An app creates the clients it needs with the address of the API and the shared configuration. After that it calls the methods like any other TypeScript function and gets typed objects back. The app does not build URLs, parse JSON or write its own types for the data.

\section{When the Client Is Generated Again}
We generate the client again every time the contract changes. That is when an endpoint is added, a route changes, a field is added or changed, or an action gets a new status code. It takes three steps:
\begin{enumerate}
  \item Start the API locally with docker compose, so it serves the OpenAPI document.
  \item Run npm run generate:api-client from the clients folder.
  \item Commit the generated file in the same pull request as the change in the API.
\end{enumerate}

We commit the generated file because the CI pipeline cannot generate it. It would need a running API inside the pipeline to do that. Instead CI type checks both apps against the committed client.

The weak point is the three steps. If a developer changes the API and forgets to generate the client again, CI does not notice, and the client stays out of date until somebody does it. The development guidelines in appendix 10.1 include the command, and a pull request that changes the API should also change the generated file.

\section{Breaking Changes}
A change is breaking when an app that worked before stops working. That happens when an endpoint or a field is removed or renamed, when a field changes type, or when an action gets a new name, because the action name becomes the method name in the client. A new endpoint or a new optional field is not breaking.

Because the client is generated, we see a breaking change right away. When we generate the client again, the TypeScript compiler fails everywhere in the apps that uses something that was removed or changed. The type check in CI then fails until the apps are fixed. So the change in the API, the new client and the fixes in the apps go into the same pull request. Our commits follow Conventional Commits (appendix 10.1), so we mark the commit as breaking with an exclamation mark after the type, for example feat(api)!, and a line that explains what breaks.

We do not version the API, because both apps are in the same repository as the API and are changed together. There is one exception. The customer app is installed on the user's phone, so an old version can still be in use after we deploy a new API. A breaking change can then stop old versions of the app from working until the user updates. We accept that in a prototype. A real product would need a versioned API, or a period where both the old and the new endpoint exist.

\section{Why One Source for the Contract}
ASP.NET Core builds the OpenAPI document from the C\# code, and NSwag generates the client from the document. That means we only write the contract once, in C\#. There are no hand-written types in the apps that can drift away from the API, and if something does not match, the compiler tells us before a user finds out. The downside is the extra generation step, and that we can forget it.

\printbibliography

\end{document}
